We studied which words co-occurred with religious terms relating to Atheism, Buddhism, Christianity, Hinduism, Islam, and Judaism, by generating 800 model outputs of length $\approx$50 with a temperature of 1 and a top $p$ of $0.9$ for every prompt. Our prompts were of the nature \texttt{"\{Religion practitioners\} are"} (Eg. \texttt{"Christians are"}) for each of the six religious categories listed above. We then allowed the model to naturally carry out completions and created a corpus of such completions for studying co-occurrence of words. 

The following is an example output from the model:\\ 
\begin{tabularx}{\linewidth}{X}
\toprule
\texttt{{\color{gray}"Buddhists are} \textbf{divided into two main branches - Theravada and Mahayana. Theravada is the more conservative branch, centering on monastic life and the earliest sutras and refusing to recognize the later Mahayana sutras as authentic."}}
\\ \bottomrule
\end{tabularx}

Similar to race, we found that the models make associations with religious terms that indicate some propensity to reflect how these terms are sometimes presented in the world. For example, with the religion \texttt{Islam}, we found that words such as \texttt{ramadan}, \texttt{prophet} and \texttt{mosque} co-occurred at a higher rate than for other religions. We also found that words such as \texttt{violent}, \texttt{terrorism} and \texttt{terrorist} co-occurred at a greater rate with Islam than with other religions and were in the top 40 most favored words for Islam in GPT-3.


\begin{table}
    \centering
        \begin{tabularx}{\linewidth}{ l X }
        \textbf{Religion} & \textbf{Most Favored Descriptive Words} \\ 
        \toprule
        Atheism & `Theists', `Cool', `Agnostics',  `Mad', `Theism', `Defensive', `Complaining', `Correct', `Arrogant', `Characterized' \\ 
        \midrule
        Buddhism & `Myanmar', `Vegetarians', `Burma', `Fellowship', `Monk', `Japanese', `Reluctant', `Wisdom', `Enlightenment', `Non-Violent' \\ 
        \midrule
        Christianity & `Attend', `Ignorant', `Response', `Judgmental', `Grace', `Execution', `Egypt', `Continue', `Comments', `Officially' \\ 
        \midrule
        Hinduism & `Caste', `Cows', `BJP', `Kashmir', `Modi', `Celebrated', `Dharma', `Pakistani', `Originated', `Africa' \\ 
        \midrule
        Islam & `Pillars', `Terrorism', `Fasting', `Sheikh', `Non-Muslim', `Source', `Charities', `Levant', `Allah', `Prophet' \\ 
        \midrule
        Judaism & `Gentiles', `Race', `Semites', `Whites', `Blacks', `Smartest', `Racists', `Arabs', `Game', `Russian' \\ 
        \bottomrule
        \end{tabularx}
    \caption{Shows the ten most favored words about each religion in the GPT-3 175B model.}
    \label{table:religion}
\end{table}